\documentclass[11pt]{article}
\usepackage[a4paper,margin=25mm]{geometry}
\usepackage[T1]{fontenc}
\usepackage[utf8]{inputenc}
\usepackage{lmodern}
\usepackage{graphicx}
\usepackage{booktabs}
\usepackage{amsmath}
\usepackage{xspace}
\usepackage{hyperref}
\graphicspath{{figs/}}
\newcommand{\code}[1]{\texttt{#1}}
\newcommand{\us}{\textmu s\xspace}
\newcommand{\Na}{\ensuremath{{}^{22}\mathrm{Na}}\xspace}
\newcommand{\Vov}{\ensuremath{V_{\rm ov}}\xspace}
\newcommand{\Vbr}{\ensuremath{V_{\rm br}}\xspace}
\newcommand{\csitl}{CsI(Tl)\xspace}

\title{CryspLight full simulation: design\\[2mm]
\large From the crystal to the ASIC words, for the comparison with the bench data}
\author{J.\,J.~Gómez-Cadenas}
\date{16 October 2026 \\[2mm] \small Version 0.1 --- written before the code}

\begin{document}
\maketitle

\begin{abstract}
The pedagogical simulation (git tag \code{pedsim}) transports scintillation
light in an ideal detector: a uniform $8\times8$ array covering the whole
readout face, one wavelength, one photon detection efficiency, photoelectrons
counted exactly. This document specifies the full simulation that produces
what the bench measures: the real SiPM layout with its gaps, the grease and the
window, the emission spectra, the SiPM as a device with gain, crosstalk and
dark counts, and the Citiroc (later the Radioroc2) chain down to the ADC
words. Every addition is switched on by a configuration key and switched off
by its absence, so the pedagogical results are reproduced bit for bit. The
measured data and the emulated data are written in the same HDF5 layout, and
one analysis reads both.
\end{abstract}

\tableofcontents

\section{Purpose and the measurements to reproduce}
\label{sec:purpose}

Three configurations, all with the same crystal size per material, the same
SiPM array, Teflon wrapping and grease coupling, read out by two Citiroc~1A
chips (32 channels each) and later by one Radioroc2 (64 channels). The ASICs
sit at ambient temperature in every case; in the cold configuration the
crystal and the SiPMs are in the cryostat and the signals come out through
cables and feedthroughs.

\begin{table}[htb]
\centering
\begin{tabular}{l c c c}
\toprule
 & BGO ambient & \csitl ambient & BGO cold \\
\midrule
crystal ($x \times y \times z$, mm) & $48\times48\times22.4$ & $48\times48\times37.2$ & $48\times48\times22.4$ \\
thickness in radiation lengths & 2 ($X_0 = 11.2$~mm) & 2 ($X_0 = 18.6$~mm) & 2 \\
temperature & laboratory (about 295~K) & laboratory & 196~K \\
light yield (ph/MeV) & 7000 to 8200$^{a}$ & 54\,000 & 14\,000 to 16\,000$^{a}$ \\
decay time & 300~ns ($+$ 60~ns minor) & 730~ns (64\%), 3340~ns (36\%) & 2.7~\us$^{b}$ \\
emission peak & 480~nm & 550~nm & 480~nm \\
refractive index & 2.15 & 1.79 & 2.15 \\
intrinsic resolution at 511~keV (FWHM) & 3 to 4\%$^{c}$ & 6\% & 3 to 4\%$^{c}$ \\
source & \multicolumn{3}{c}{\Na, on axis, 30~mm from the face opposite the SiPMs} \\
readout & \multicolumn{3}{c}{Citiroc 1A $\times 2$, then Radioroc2} \\
\bottomrule
\end{tabular}
\caption{The three configurations. $^{a}$ Parameter: the ambient yield from the
Gironnet quenching curve relative to the Ding calibration at 195~K; the cold
yield between the Ding chain (14k) and the value implied by the longer decay
time measured on our crystal. The warm-to-cold ratio of the measured
photoelectrons fixes it. $^{b}$ Measured on this crystal in the Antarctica runs
($2.72 \pm 0.01$~\us at T3 = 196~K); replaces the 1768~ns of the Ding
calibration for this crystal only. $^{c}$ Parameter; BGO is nearly
proportional, the value is to be taken from the data once the statistical and
noise terms are subtracted.}
\label{tab:configs}
\end{table}

The simulation has three stages with a file between each pair, so that a
change downstream never requires rerunning upstream:

\begin{enumerate}
\item \textbf{Optical stage} (Julia, the existing pipeline extended). Gammas
  from the source, energy deposits, scintillation photons transported on the
  GPU to the SiPMs. Output: the truth already written by \code{pedsim}, plus
  one record per detected photoelectron (channel, time, wavelength).
\item \textbf{Device and front-end stage} (Julia, new module). SiPM response
  and the electronics chain applied to the photoelectron records. Output: the
  words the DAQ would write, per event and channel, in the common HDF5 layout.
\item \textbf{Comparison stage} (Python, \code{fullsim\_analysis/}). The same
  code on the emulated and the measured files: gains, pedestals, spectra,
  resolution, maps.
\end{enumerate}

\section{Principles}
\label{sec:principles}

\begin{itemize}
\item \textbf{Additive.} No existing behaviour changes. Each new feature is a
  configuration block; absent, the code follows the \code{pedsim} path with the
  same random draws. A frozen-output test runs a small pedagogical
  configuration and compares its output with stored checksums.
\item \textbf{Kernel invariant kept.} Every physics addition is written in both
  \code{src/transport.jl} (reference) and \code{src/kernel.jl} (GPU), and the
  parity test proves them identical on the CPU backend.
\item \textbf{Measured where measurable.} Noise, gains, breakdown voltages,
  thresholds and hold timing are inputs read from the run, never datasheet
  defaults, once the run exists. Until then the defaults are stated here with
  their provenance.
\item \textbf{One layout for data and emulation.} Section~\ref{sec:formats}.
\end{itemize}

\section{Geometry}
\label{sec:geometry}

\subsection{Crystal and wrapping}

A rectangular box, $48 \times 48$~mm$^2$ transverse, 22.4~mm (BGO) or
37.2~mm (\csitl) thick, polished on all six faces. The polish is the measured
one of the pedagogical calibration: a micro-facet Gaussian with
$\sigma_\alpha = 1.3^\circ$. Five faces are hand-wrapped in several layers of
PTFE tape: a thin air gap between crystal and tape (the UNIFIED
``back-painted'' finish, the \code{pedsim} default) and a Lambertian
reflector of $R = 0.99$ behind it (Janecek and Moses; saturated for three or
more layers). Nothing here is new.

\subsection{Coupling and window}

The readout face is coupled through optical grease, $n = 1.42$, to the SiPM
tiles, whose photosensitive surface is covered by 0.15~mm of silicone resin
with $n = 1.57$ (datasheet). The photon crosses two interfaces, crystal to
grease and grease to resin, each treated with Fresnel reflection and
transmission; the first dominates (BGO 2.15 to 1.42; \csitl 1.79 to 1.42),
the second reflects 0.3\% at normal incidence. The resin-to-silicon step is
not simulated: the datasheet PDE is measured from air through the resin and
already contains it. The PDE is applied once the photon is in the resin above
an active area, as a function of wavelength and overvoltage
(Section~\ref{sec:sipm}) and independent of angle; this angular assumption
is the one unverified element of the coupling model and is listed in
Section~\ref{sec:open}. Reflected photons return to the crystal and continue.

\subsection{SiPM layout}
\label{sec:layout}

Four Hamamatsu S14161-6050HS-04 tiles of $4\times4$ channels, 6.0~mm active
width on a 6.2~mm pitch inside a tile, 0.4~mm between the active areas of
neighbouring tiles. The 64 centres, with the origin at the centre of the
matrix, are at $\pm 3.2, \pm 9.4, \pm 15.6, \pm 21.8$~mm in each coordinate.
Consequences for the model:

\begin{itemize}
\item the matrix is 49.6~mm wide: the outer channels overhang the crystal by
  0.8~mm and only 5.2~mm of their width sees the crystal;
\item inside the $48 \times 48$~mm$^2$ footprint the active fraction is
  $0.934$: 0.2~mm gaps between channels, the 0.4~mm seam through the centre
  in both directions, and the overhang;
\item a photon transmitted into the resin above a dead strip meets the
  package surface between the channels. Its reflectivity is unknown; the
  default is diffuse with $R_{\rm dead} = 0.5$, and
  Section~\ref{sec:comparison} says how the data constrain it.
\end{itemize}

The layout enters the configuration as a list of active rectangles with a
channel number each. The \code{pedsim} path (pitch 6.0, no gaps) is the
default when the block is absent. The channel numbering follows the DAQ
cabling and is part of the run description (Section~\ref{sec:formats}).

\section{Source and gamma stage}
\label{sec:gamma}

A \Na point source on the crystal axis, 30~mm from the centre of the face
opposite the SiPMs. Per decay: a 1274.5~keV gamma (99.9\%), and with 90.3\%
probability a positron whose annihilation gives two back-to-back 511~keV
gammas. Directions are isotropic; only the gammas whose line meets the front
face are transported (13\% of the sphere for the 511~keV pair, counted once
since the two partners cannot both reach the crystal; 13\% for the 1275~keV
gamma). The two lines of one decay are kept together so that true
coincidence summing (the 1786~keV sum peak) and the 1275~keV Compton
continuum under the 511~keV line appear with their real weights. Incidence
angles reach $49^\circ$ at the corners.

The gamma transport is the vendored PTCrysp core (see
\code{src/gamma/VENDORED.md}): photoelectric absorption, Compton scattering,
pair production above 1.022~MeV treated as local absorption (at 1275~keV it
is below 1\% of the cross section), tables to 10~MeV. Not modelled, and
listed in Section~\ref{sec:open}: K-shell fluorescence and its escape (the
Bi K X-ray at 77~keV gives an escape peak at 434~keV for interactions within
a fraction of a millimetre of a surface), Doppler broadening, scattering in
the source holder and the cryostat walls.

The \code{pedsim} source (511~keV, uniform over the front face, normal
incidence) remains the default; the point source is a \code{[source]} block.

\section{Scintillation}
\label{sec:scintillation}

Each deposit emits a Poisson number of photons with mean yield times energy
(Table~\ref{tab:configs}); the intrinsic resolution of the crystal, which the
optics cannot produce, is applied as a Gaussian smearing of the yield per
event with the FWHM of the table, a parameter per material. Emission times
follow the decay profile: a single exponential for BGO, two components for
\csitl, sampled per photon, added to the deposit time. Wavelengths are drawn
from the tabulated emission spectrum of each material; the refractive index,
the Rayleigh length ($\propto \lambda^{-4}$ from the calibrated value at the
peak) and the PDE then depend on the drawn wavelength. Bulk absorption keeps
the calibrated scalar value. The scalar \code{pedsim} path (one wavelength,
one decay time) is the default.

The cold BGO yield and decay time are a new material file; the pedagogical
\code{data/bgo.toml} (Ding calibration, 1768~ns) is untouched.

\section{Optical transport}
\label{sec:transport}

The transport is the \code{pedsim} kernel: Philox counter-based random
numbers per photon, the box geometry, UNIFIED surfaces with the air gap and
the specular lobe, Rayleigh scattering, bulk absorption, Fresnel at the
readout face, exact photon conservation across the terminal statuses. Three
additions, each gated:

\begin{enumerate}
\item the active-rectangle layout and the dead-area reflection
  (Section~\ref{sec:layout}),
\item the second interface and the resin layer (the photon crosses it with no
  lateral displacement worth tracking: 0.15~mm over a 6~mm pitch),
\item per-photon wavelength with the tabulated $n(\lambda)$,
  $\Lambda_{\rm Rayleigh}(\lambda)$ and PDE($\lambda$).
\end{enumerate}

Metal constraints: tables travel packed in one Float32 buffer, flags and
scalars in one isbits struct, no dynamic tuple indexing. New random draws
occur only inside enabled branches.

Output of the stage, per event: the \code{pedsim} truth (deposits, first and
second interaction, energies, map, first-photoelectron times) unchanged, and
a ragged list of detected photoelectrons: channel (UInt8), time (Float32,
ns), wavelength (UInt16, nm). At 7 bytes per photoelectron a 511~keV event
is about 70~kB in \csitl and 15~kB in BGO; $10^5$ events per configuration fit
in a few gigabytes and take minutes on the GPU. The pedagogical
\code{events.h5} is not changed; the records go to \code{pes.h5} in the same
output directory.

\section{SiPM response}
\label{sec:sipm}

The device is the S14160-6050HS element (datasheet KAPD1064E, November
2024; the S14161 tiles carry the same element). The model, per channel:

\begin{description}
\item[Breakdown voltage.] $\Vbr(T) = 38.0~\mathrm{V} + 0.034~\mathrm{V/K}\,(T - 298~\mathrm{K})$,
  typical; 34.5~V at 196~K, measured 34.3~V on the Antarctica channel~49.
  Per-channel spread 0.1~V typical, 0.2~V maximum within a tile. The
  simulation takes $\Vbr$ per channel from the run (measured as in
  \code{overvoltage\_recipe.tex}) or draws it from the spread.
\item[Overvoltage.] $\Vov = V_{\rm bias} - \Vbr$ per channel, after the
  input-DAC trim if used.
\item[Gain.] $G = 0.93 \times 10^6\,\Vov/\mathrm{V}$ (linear, zero at
  breakdown); $2.5 \times 10^6$ at 2.7~V.
\item[PDE.] $\mathrm{PDE}(\lambda, \Vov) = \mathrm{PDE}_{450}(\Vov) \times s(\lambda)$,
  with $\mathrm{PDE}_{450}$ from the datasheet curve (37\% at 1.7~V, 50\% at
  2.7~V, 54\% at 3.2~V, flat above) and $s(\lambda)$ the normalised spectral
  curve (digitised from the datasheet figure).
\item[Crosstalk.] Probability $p_{\rm xt} = 0.05\,\Vov/\mathrm{V} - 0.06$
  (2.5\% at 1.6~V, 7\% at 2.7~V, 13\% at 3.7~V), applied as a geometric
  number of extra, simultaneous avalanches per primary.
\item[Afterpulsing.] Not quoted by Hamamatsu for this series beyond ``low'';
  default 2\% with a 100~ns time constant, parameter.
\item[Dark counts.] Rate per device $R(T, \Vov)$, from the typical dark
  current at 25~$^\circ$C and 2.7~V ($2.5~\mu$A, i.e.\ $6 \times 10^6$
  avalanches per second after crosstalk), scaled with $\Vov$ like the PDE and
  with temperature by a factor 2 per 9~K; anchored cold by the measured 8 to
  15~kHz per channel at 196~K and 6.5~V. Dark avalanches are added to the
  photoelectron train over the whole acquisition window, with crosstalk.
\item[Single-photoelectron pulse.] The current pulse of one avalanche into the
  front-end input impedance; the shape measured in the Antarctica runs into
  100~$\Omega$ (rise 16~ns, 24~ns FWHM, 50~ns decay) is the default template,
  with the amplitude scaled by the gain. The Citiroc input (50~$\Omega$ and
  100~nF in series, into a voltage preamplifier with 15~pF input capacitance)
  sees a different shape; it is a parameter until measured.
\item[Cell occupancy.] 14\,331 cells per device; the probability that a
  photoelectron lands on a cell still recovering is tracked with a 50~ns
  recovery. Only the brightest \csitl channel approaches a percent effect.
\end{description}

The stage output is, per channel, a train of avalanche times with their
charge (one or more cells, after crosstalk), which is the input to the
electronics.

\section{Electronics}
\label{sec:electronics}

\subsection{Citiroc 1A}

Two chips read the 64 channels (32 each, channels assigned by the cabling in
the run description). Per channel, the chain of the datasheet (DS7004 v2.5)
as emulated:

\begin{enumerate}
\item \textbf{Input.} The avalanche train becomes a voltage at the preamplifier
  input through the termination; the input DAC (8 bits, 2.5 or 4.5~V range)
  only shifts the SiPM bias and enters through \Vov.
\item \textbf{Preamplifiers.} High gain (15~pF input capacitance, gain 10 to
  600 by feedback capacitor) and low gain (1.5~pF, gain 1 to 60) in parallel,
  per-channel settings; a gain-spread parameter per channel.
\item \textbf{Slow shapers.} CR-RC$^2$ with peaking time $t_p$ from 12.5 to
  87.5~ns in 12.5~ns steps, one value per gain path common to the chip. The
  shaped signal is the convolution of the input with the normalised impulse
  response $(t/\theta)^2 e^{-t/\theta}$, $t_p = 2\theta$. The order of the
  shaper is what the datasheet states; its exact transfer function is a
  parameter to be checked against a single-photoelectron waveform if one can
  be probed.
\item \textbf{Noise.} Added after the shaper as a Gaussian series with the
  spectrum measured on pedestal or control events of the run; default white
  with rms 0.1 photoelectron on the high-gain path (the datasheet
  signal-to-noise of 10 for one photoelectron).
\item \textbf{Trigger.} Fast shaper (15~ns peaking) on the high-gain path,
  two discriminators with chip-wide 10-bit thresholds and 4-bit per-channel
  trims: the time discriminator gives the per-channel trigger output, the
  charge discriminator feeds the hit register and the chip OR. The emulation
  provides, per channel, the first threshold crossing of the fast-shaped
  signal; the acquisition trigger is a configurable function of them (OR of
  one chip, OR of both, majority $N$ within a coincidence window), built by
  the FPGA in the real system.
\item \textbf{Peak detector and hold.} Armed at the trigger, it stores the
  maximum of each shaped signal (high and low gain separately, all 32
  channels of a chip at once) until the hold, a fixed delay after the trigger
  (parameter; up to 500~ns in our DAQ). Each of the 64 channels stores its
  own maximum, including the maximum of its noise when it holds little
  light: the positive bias this gives to dim channels is in the emulation by
  construction. Track-and-hold mode is available as an alternative.
\item \textbf{Readout and ADC.} The stored values are multiplexed at up to
  5~MHz and digitised externally; the emulation applies the ADC gain,
  offset, 12-bit quantisation (parameter) and the saturation of the chain,
  400~pC of input charge or the ADC full scale, whichever comes first, per
  path. Readout dead time is recorded but not simulated as an efficiency:
  the source rates are low.
\end{enumerate}

With $t_p = 87.5$~ns the shaper integrates about $e\,t_p = 240$~ns of light.
The stored fraction of the photoelectrons is then about 0.4 for BGO at
ambient, 0.24 for \csitl (the 730~ns component carries 64\%), 0.09 for BGO at
2.7~\us, and the statistical floor of the resolution follows from the
photoelectrons inside that fraction; the emulation computes it with the 64
independent channels, the noise and the dark pulses.

\subsection{Radioroc2}

One chip (datasheet v2.0, January 2025). Same structure, with the differences
that matter: 64 channels; per-channel shaping time from 20 to 300~ns, or 120~ns
to 1.8~\us in the slow mode (the natural choice for \csitl and cold BGO);
charge preamplifiers with per-channel gains (high 5 to 80, low 0.5 to 8);
two time discriminators (thresholds T1 and T2) and a charge discriminator on
the low-gain shaper; the hold generated on chip by a delay cell (2~ns to
2.6~\us) or externally; saturation at 2000 photoelectrons times $10^6$ gain
(about 320~pC); the option of routing each channel's shaped waveform to its
own output pin. The emulation is the same code with a different parameter
block; the shaper order is assumed CR-RC$^2$ until measured.

\subsection{Not planned}

A waveform-digitiser variant (window integration with measured $1/f$ noise,
majority trigger on dark pulses) was considered for the Antarctica runs and
dropped: those runs are not to be simulated.

\section{File formats}
\label{sec:formats}

\subsection{Photoelectron records (simulation only)}

\code{output/<tag>/pes.h5}: datasets \code{offset} (n\_events+1, Int64),
\code{channel} (UInt8), \code{t\_ns} (Float32), \code{lambda\_nm} (UInt16),
with the event index aligned with \code{events.h5}; attributes copy the
configuration and the git revision.

\subsection{Event words: one layout for measured and emulated data}

The measured Citiroc files are converted, as the first step of the analysis,
to this layout; the front-end stage writes it directly. HDF5, one file per
run:

\begin{description}
\item[\code{/run} attributes] date, crystal, material, temperature readings and
  their placement, bias per chip, measured \Vbr per channel, overvoltage per
  channel, thresholds, hold delay, shaping times, preamplifier gains, ADC
  description, DAQ trigger definition, git revision of the converter or the
  emulator; for emulated runs, the full simulation configuration and the
  source of every parameter.
\item[\code{/channels}] \code{x\_mm}, \code{y\_mm} (64), \code{tile} (64),
  \code{chip}, \code{asic\_channel} (64), \code{status} (64: live, dead,
  noisy), the mapping from DAQ channel to the geometry of
  Section~\ref{sec:layout}.
\item[\code{/calib}] \code{spe\_gain\_hg}, \code{spe\_gain\_lg} (64, ADC per
  photoelectron, with errors), \code{pedestal\_mean}, \code{pedestal\_rms} (64,
  per path), \code{dark\_rate} (64, Hz), measured on the run or on its
  calibration runs; for emulated runs, the true values used.
\item[\code{/events}] \code{hg} (n, 64, UInt16), \code{lg} (n, 64, UInt16),
  \code{hit} (n, 64, Bool: the hit register), \code{t\_trig} (n), per-event
  flags (saturation per path, pile-up); and, for emulated runs only, a
  \code{/truth} group with the \code{events.h5} truth (interaction positions,
  energies, total photoelectrons, the integrated map) so that any observable
  can be resolved against the truth.
\end{description}

Arrays are event-major, so \code{h5py} reads $(n, 64)$ directly. The pedagogical
\code{events.h5} layout is not modified.

\section{Comparison observables}
\label{sec:comparison}

The same scripts on both kinds of file, in this order; each line names what
the observable constrains.

\begin{enumerate}
\item \textbf{Single-photoelectron spectra} from dark pulses, per channel: the
  gain table and its spread; against bias, the breakdown voltage per channel.
  Constrains \Vov and the gain model; it is also the thermometer of the
  SiPMs.
\item \textbf{Pedestals and control charges}: the noise term per channel and
  on the sum. Measured input to the emulation, not a prediction.
\item \textbf{Dark rate per channel} against temperature and bias.
\item \textbf{Total charge spectrum} (equalised sum of the 64 channels, high
  and low gain): the 511 and 1275~keV lines, the sum peak, the Compton
  continua, the backscatter region. The ratio of the two photopeaks (2.5 if
  linear) tests the dynamic range and the trigger; the line widths and their
  decomposition (noise from the pedestals, statistics from the stored
  fraction, intrinsic from what remains) test the front-end model.
\item \textbf{Light yield}: photoelectrons at 511~keV, converted through the
  simulated collection efficiency and PDE to photons per MeV; the
  warm-to-cold BGO ratio against the quenching curve.
\item \textbf{Brightest-channel distribution}: its tail and the saturation kink
  in \csitl; the ratio of brightest channel to total against the simulated
  light maps tests the optics near the readout face.
\item \textbf{Occupancy and mean-charge maps} per channel: the edge channels
  (5.2~mm under the crystal), the central seam, the tile boundaries. The
  dead-area reflectivity $R_{\rm dead}$ and the surface roughness are the
  free parameters these maps constrain.
\item \textbf{Seed-channel photopeak map}: the photopeak position of the events
  grouped by their brightest channel, in units of the median. In a healthy
  array this is a map of the collection non-uniformity (edges, seam), to be
  reproduced by the simulation, not corrected away.
\item \textbf{Overvoltage dependence}, when a bias scan exists: light,
  crosstalk excess, resolution, saturated fraction against \Vov, against the
  curves of Section~\ref{sec:sipm}.
\end{enumerate}

\section{Validation}
\label{sec:validation}

\begin{itemize}
\item Frozen-output test: a small \code{pedsim} configuration reproduces stored
  checksums of its maps and truth after every change.
\item Parity test (existing): reference and kernel identical on the CPU
  backend, extended to every gated branch switched on.
\item Layout test: the active-rectangle layout with the \code{pedsim} rectangles
  gives the \code{pedsim} maps; with gaps, the dead fraction of detected
  photons at $R_{\rm dead} = 0$ equals the geometric dead fraction for a
  uniform illumination.
\item Window test: with $n_{\rm grease} = n_{\rm resin}$ the second interface
  reflects nothing; with the resin index equal to the crystal index and no
  grease the result equals the contact case.
\item Spectral test: a delta emission spectrum reproduces the scalar run.
\item Photoelectron records test: the records re-binned in 100~ns give the
  \code{pedsim} maps exactly.
\item Source test: the 511~keV rate on the face equals the solid-angle fraction;
  the incidence-angle distribution matches the analytic one.
\item Front-end unit tests: a delta input gives the analytic CR-RC$^2$ peak at
  $t_p$; a Poisson train with a long decay gives the stored fraction of
  Section~\ref{sec:electronics} within statistics; the noise-only peak
  detector gives the expected maximum-of-Gaussian bias; the saturation clips
  at the configured charge.
\item End-to-end analytic check: the energy resolution of a configuration from
  the emulation against the quadrature of its stated terms.
\end{itemize}

\section{Parameters and provenance}
\label{sec:parameters}

\begin{table}[htb]
\centering
\small
\begin{tabular}{l l l l}
\toprule
parameter & value & source & status \\
\midrule
crystal sizes & Table~\ref{tab:configs} & setup & fixed \\
polish $\sigma_\alpha$ & $1.3^\circ$ & pedsim calibration & fixed \\
Teflon & Lambertian, $R = 0.99$, air gap & Janecek \& Moses; pedsim & fixed \\
grease index & 1.42 & setup & fixed \\
resin index, thickness & 1.57, 0.15~mm & datasheet & fixed \\
SiPM layout & centres at $\pm3.2, \pm9.4, \pm15.6, \pm21.8$~mm & setup & fixed \\
dead-area reflectivity & 0.5, diffuse & assumption & to fit on maps \\
BGO yield ambient & 7000 to 8200 ph/MeV & Gironnet curve, Moszynski & to fit on data \\
BGO yield 196~K & 14\,000 to 16\,000 ph/MeV & Ding chain; Antarctica $\tau$ & to fit on data \\
BGO decay 196~K & 2.7~\us & Antarctica runs & fixed \\
BGO decay ambient & 300~ns & literature & fixed \\
\csitl yield, decay & 54\,000; 730/3340~ns, 64/36\% & pedsim calibration; literature & fixed \\
intrinsic resolution & 6\% \csitl; 3 to 4\% BGO & CRYSP Citiroc data; assumption & BGO to fit \\
emission spectra & tabulated & literature & to digitise \\
\Vbr(T) & $38~\mathrm{V} + 34~\mathrm{mV/K}\,(T-298~\mathrm{K})$ & datasheet; Antarctica & measured per run \\
gain, PDE, crosstalk vs \Vov & Section~\ref{sec:sipm} & datasheet curves & fixed \\
dark rate & $6$~MHz at 25~$^\circ$C, 2.7~V; 10~kHz at 196~K, 6.5~V & datasheet; Antarctica & measured per run \\
afterpulsing & 2\%, 100~ns & assumption & parameter \\
SPE pulse shape & 16/24/50~ns into 100~$\Omega$ & Antarctica & parameter for Citiroc input \\
Citiroc shaper & CR-RC$^2$, $t_p = 87.5$~ns & datasheet & setting of the run \\
Citiroc noise & 0.1 pe rms on HG & datasheet SNR & measured per run \\
hold delay & up to 500~ns & our DAQ & setting of the run \\
ADC & 12 bit & assumption & setting of the run \\
source distance & 30~mm, on axis, far face & setup & fixed \\
\bottomrule
\end{tabular}
\caption{Every parameter of the full simulation with its value, where it comes
from, and whether it is fixed, measured on the run, or to be fitted on the
comparison.}
\label{tab:params}
\end{table}

\section{Open items}
\label{sec:open}

\begin{itemize}
\item Angular dependence of the PDE through the resin when illuminated from
  grease rather than air (assumed none).
\item Reflectivity and specularity of the package surface between channels.
\item BGO light yield at ambient and at 196~K for this crystal; whether the
  2.7~\us decay implies an effective temperature below T3.
\item BGO intrinsic resolution.
\item Citiroc: shaper transfer function beyond ``CR-RC$^2$'', noise spectrum,
  peak-detector droop over the readout, pulse shape at the preamplifier
  input with a 2~nF device.
\item Radioroc2: shaper order, noise, 100~$\Omega$ input mode with 2~nF, DAQ
  board and ADC.
\item Gamma stage: K-shell fluorescence and escape, Doppler broadening,
  scattering in the source holder and cryostat.
\item Temperature of the SiPMs in the cold run relative to the crystal
  (resolved operationally by the measured \Vbr).
\end{itemize}

\end{document}
